% 方法命名：MOEA/D–ALNS为原始参数配置的ALNS基线；MOEA/D–ALNS-T为经开发集参数优化的ALNS基线；MOEA/D–DPR为本文采用学习模型指导破坏操作的MOEA/D方法。三种方法均使用MOEA/D外层框架；ALNS-T中的T表示参数优化，不表示其余方法不含自适应操作选择。MOEA/D–DPR为本地学习增强实现的论文简称，不表述为原作者DPR的原样复现。
% Requires booktabs; Chinese typesetting via ctex or equivalent.
\begin{table}[htbp]
\centering
\caption{三种方法的HV与IGD+比较（均值$\pm$标准差）}
\label{tab:three-method-quality}
\begin{tabular}{llcc}
\toprule
客户数 & 方法 & HV$\uparrow$ & IGD+$\downarrow$ \\
\midrule
50 & MOEA/D–ALNS & $0.92484 \pm 0.00955$ & $0.06942 \pm 0.00531$ \\
50 & MOEA/D–ALNS-T & $0.93992 \pm 0.00831$ & $0.06140 \pm 0.00490$ \\
50 & MOEA/D–DPR & $\mathbf{0.96700 \pm 0.01140}$ & $\mathbf{0.05387 \pm 0.00694}$ \\
\midrule
150 & MOEA/D–ALNS & $0.51381 \pm 0.01802$ & $0.16656 \pm 0.01234$ \\
150 & MOEA/D–ALNS-T & $0.54974 \pm 0.01178$ & $0.14353 \pm 0.00891$ \\
150 & MOEA/D–DPR & $\mathbf{0.59756 \pm 0.01581}$ & $\mathbf{0.11818 \pm 0.00957}$ \\
\midrule
250 & MOEA/D–ALNS & $0.49186 \pm 0.02231$ & $0.20806 \pm 0.01211$ \\
250 & MOEA/D–ALNS-T & $0.54842 \pm 0.02061$ & $0.17917 \pm 0.00873$ \\
250 & MOEA/D–DPR & $\mathbf{0.59366 \pm 0.03510}$ & $\mathbf{0.13083 \pm 0.01686}$ \\
\midrule
总体 & MOEA/D–ALNS & $0.64350 \pm 0.00965$ & $0.14801 \pm 0.00577$ \\
总体 & MOEA/D–ALNS-T & $0.67936 \pm 0.00774$ & $0.12804 \pm 0.00525$ \\
总体 & MOEA/D–DPR & $\mathbf{0.71941 \pm 0.01098}$ & $\mathbf{0.10096 \pm 0.00560}$ \\
\bottomrule
\end{tabular}
\par\smallskip
\begin{minipage}{0.98\linewidth}\footnotesize
注：每种规模包含15个实例，每实例20个搜索种子。HV越大越好，IGD+越小越好；加粗表示同规模最优均值。各实例独立构建三方法×20seed共同经验非支配参考，每个seed档案独立计分后对实例等权汇总。±为20个seed各自跨实例宏观均值的样本标准差，不是900次运行直接混合的SD，也不是独立训练重复的SD。
\end{minipage}
\end{table}

\begin{table}[htbp]
\centering
\caption{MOEA/D–DPR与两种ALNS的总体配对Wilcoxon符号秩检验}
\label{tab:three-method-test}
\begin{tabular}{llcrrr}
\toprule
对照 & 指标 & 胜/平/负 & $W$ & $p$ & $p_{\mathrm{Holm}}$ \\
\midrule
MOEA/D–ALNS & HV & 40/0/5 & 84 & 6.6549e-08 & 1.9965e-07 \\
MOEA/D–ALNS & IGD+ & 39/0/6 & 80 & 4.5573e-08 & 1.8229e-07 \\
MOEA/D–ALNS-T & HV & 30/0/15 & 276 & 0.0057111 & 0.0057111 \\
MOEA/D–ALNS-T & IGD+ & 32/0/13 & 238 & 0.0012307 & 0.0024613 \\
\bottomrule
\end{tabular}
\par\smallskip
\begin{minipage}{0.98\linewidth}\footnotesize
注：以45个实例的20seed指标均值为配对单位，采用双侧精确条件符号秩检验。胜/平/负均以MOEA/D–DPR为参照。Holm校正覆盖三方法两两比较×HV/IGD+共6项检验，表中仅展示MOEA/D–DPR相关四项。本表为用户指定三方法集合的事后补充分析，原计划检验保留不变。
\end{minipage}
\end{table}
